% SOURCES: draft/Features/features.md (Tier-1 capabilities); docs/0.2.0v/plan.md
%          (§11 Feature Checklist, §12 Additional In-Scope, §5 Architecture);
%          README.md; docs/1.0.0v/supervisor-report.md
% GUIDANCE: Split into functional + non-functional requirements, then present the
% primary USE CASES (the marking scheme asks for these). A use-case diagram goes
% in ch. 7 (Analysis); here, write the use-case descriptions.

\chapter{Requirements Gathering}
\label{ch:requirements}

\section{Functional Requirements}
% DRAFT (review): derived from the Tier-1 feature set.
\begin{enumerate}[label=\textbf{FR\arabic*.}, leftmargin=3.4em]
  \item The system shall let a user insert, select, move, reorder, group, and
        delete elements on a canvas.
  \item The system shall support primitive elements (container, text, image,
        button, link, icon, list, divider) and \emph{presets} that compose them.
  \item The system shall let a user define design tokens and bind element styles
        to them, with light/dark theming.
  \item The system shall support per-breakpoint style values
        (desktop/tablet/mobile/small) and element states (hover/focus/active).
  \item The system shall create an element from a drawn rectangle (draw-to-create).
  \item The system shall match a drawn layout to a bundled professional design
        (match-layout).
  \item The system shall validate the document and export a portable HTML/CSS/JS
        bundle as a ZIP archive.
  \item The system shall expose the document/generate/export pipeline over a
        headless MCP server.
  \item The system shall save/open projects as versioned \texttt{.dtw} files with
        crash recovery and external-change reload.
\end{enumerate}

\section{Non-Functional Requirements}
% SOURCES: plan.md §14 Performance Budgets, §15 Security; CLAUDE.md (constraints)
% DRAFT (review):
\begin{enumerate}[label=\textbf{NFR\arabic*.}, leftmargin=3.8em]
  \item \textbf{Accessibility:} export is blocked on any critical/serious
        \texttt{axe-core} violation; WCAG-aligned output \cite{wcag21}.
  \item \textbf{Determinism:} identical input yields byte-identical output.
  \item \textbf{Performance:} data-layer and export operations meet defined
        budgets (see ch.~\ref{ch:testing}).
  \item \textbf{Security:} Electron runs with context isolation, no node
        integration, sandboxing, a strict CSP, and Zod-validated IPC boundaries.
  \item \textbf{Portability:} output has no build step and no mandatory runtime
        dependencies.
\end{enumerate}

\section{Primary Use Cases}
% GUIDANCE: one short structured description per use case (actor, precondition,
% main flow, postcondition). These feed the use-case diagram in ch. 7.
% DRAFT (review) — expand the flows:

\subsection*{UC-1: Author and export a page}
\textbf{Actor:} Author. \textbf{Pre:} app open. \textbf{Flow:} open a template or
blank page → insert/edit elements → adjust tokens/theme/breakpoints → run
validation → export ZIP. \textbf{Post:} a valid, accessible bundle on disk.

\subsection*{UC-2: Draw to create an element}
\textbf{Actor:} Author. \textbf{Flow:} select a container → drag a rectangle → the
system ranks the likely element kind with a confidence hint → author accepts →
the element is placed via a normal insert operation.

\subsection*{UC-3: Match my layout}
\textbf{Actor:} Author. \textbf{Flow:} rough out a layout → open Match → the system
fingerprints the tree and ranks bundled designs with a per-dimension breakdown →
author adopts one → it hydrates as a normal document.

\subsection*{UC-4: Drive the builder headlessly (MCP)}
\textbf{Actor:} MCP client / agent. \textbf{Flow:} \texttt{create\_document} →
insert presets / set tokens / set theme → \texttt{run\_a11y\_check} →
\texttt{export\_site}, all without the GUI.

% TODO: add UC-5 (save/open project, crash recovery) and any others your marker
% expects; then draw the use-case diagram in ch. 7.
